\documentclass[a4paper, 12pt]{book} % changed from scrbook
% not using \documentclass[a4paper, twoside=semi]{scrreprt} anymore

\usepackage{../Style-Files/InfoSec}

\definecolor{bg}{RGB}{20,20,20}
\definecolor{fg}{RGB}{220,220,220}
\definecolor{keyword}{RGB}{130,170,255}
\definecolor{comment}{RGB}{120,120,120}
\definecolor{string}{RGB}{255,200,120}

\lstdefinestyle{darkbash}{
    backgroundcolor=\color{bg},
    basicstyle=\scriptsize\ttfamily\color{fg},
    keywordstyle=\color{keyword},
    commentstyle=\color{comment},
    autogobble=true,
    stringstyle=\color{string},
    breaklines=true,
    showstringspaces=false,
    frame=single
}

% \usepackage{titlesec}
% \setkomafont{chapter}{\Huge\bfseries}
% \setkomafont{section}{\Large\bfseries}
% \setkomafont{subsection}{\large\bfseries}
\usepackage{fancyhdr, lastpage}
\usepackage{lstautogobble}
\usepackage{euscript}


\makeatletter
\providecommand{\subtitle}[1]{\gdef\@subtitle{#1}}
\providecommand{\@subtitle}{}
\makeatother


\usepackage{booktabs}
\usepackage{array}
\usepackage{multirow}




\usepackage{newfloat}
%\usepackage{tocloft}

% % Define new float environment
% \DeclareFloatingEnvironment[
%   fileext=lof,
%   listname={List of Flowcharts},
%   name=Flowchart,
% ]{flowchart}


\pagestyle{fancy}
%\fancyhf{}
\lhead{}
\chead{}
%\rhead{Introduction to the Theory of Computation}
\cfoot{\thepage\ of \pageref{LastPage}}
\rfoot{Version 8.3}




\RequirePackage{etex} 


%--- The documentation for the title can be found at: https://en.wikibooks.org/wiki/LaTeX/Title_Creation
\title{Information Security \\ OpenSSL}
\author{Ahmed U. Ashraf}
\date{\today }
% \subtitle{Version~\version}

\begin{document}

% \include{glossary}


\maketitle
%--- we would like to hyperlink each item in contensts page
% \let\hypercontentsline=\contentsline
% \renewcommand{\contentsline}[4]{\hypertarget{toc.#4}{}\hypercontentsline{#1}{#2}{#3}{#4}}
%--- We print ToC: Table of Contents

\DeclareFontShape{U}{stmry}{b}{n}{<-> ssub * stmry/m/n}{}

\tableofcontents\label{toc}
%--- We print the list of figures, this might go at the end.




% \listoffigures
% %--- We print the list of tables, this might also go at the end.
% \listoftables
% %--- This is to add list of symbols as glossary
% % \printunsrtglossary[type=symbols,style=long,title={List of Symbols}]


% % List of Algorithms (- taken the idea from Discrete Algorithmic Mathematics by Stephen B.Maurer and ANthony Ralston.pdf)

% \listofalgorithms
% \addcontentsline{toc}{chapter}{List of Algorithms}

% \listofflowcharts
% %\addcontentsline{toc}{chapter}{List of Flowcharts}



\chapter{Encodings}




\verb|openssl enc -list|


\begin{table}[H]
\centering
\small
\begin{tabular}{llll}
\hline
aes-128-cbc & aes-128-cfb & aes-128-cfb1 & aes-128-cfb8 \\
aes-128-ctr & aes-128-ecb & aes-128-ofb & aes-192-cbc \\
aes-192-cfb & aes-192-cfb1 & aes-192-cfb8 & aes-192-ctr \\
aes-192-ecb & aes-192-ofb & aes-256-cbc & aes-256-cfb \\
aes-256-cfb1 & aes-256-cfb8 & aes-256-ctr & aes-256-ecb \\
aes-256-ofb & aes128 & aes128-wrap & aes128-wrap-pad \\
aes192 & aes192-wrap & aes192-wrap-pad & aes256 \\
aes256-wrap & aes256-wrap-pad & aria-128-cbc & aria-128-cfb \\
aria-128-cfb1 & aria-128-cfb8 & aria-128-ctr & aria-128-ecb \\
aria-128-ofb & aria-192-cbc & aria-192-cfb & aria-192-cfb1 \\
aria-192-cfb8 & aria-192-ctr & aria-192-ecb & aria-192-ofb \\
aria-256-cbc & aria-256-cfb & aria-256-cfb1 & aria-256-cfb8 \\
aria-256-ctr & aria-256-ecb & aria-256-ofb & aria128 \\
aria192 & aria256 & bf & bf-cbc \\
bf-cfb & bf-ecb & bf-ofb & blowfish \\
camellia-128-cbc & camellia-128-cfb & camellia-128-cfb1 & camellia-128-cfb8 \\
camellia-128-ctr & camellia-128-ecb & camellia-128-ofb & camellia-192-cbc \\
camellia-192-cfb & camellia-192-cfb1 & camellia-192-cfb8 & camellia-192-ctr \\
camellia-192-ecb & camellia-192-ofb & camellia-256-cbc & camellia-256-cfb \\
camellia-256-cfb1 & camellia-256-cfb8 & camellia-256-ctr & camellia-256-ecb \\
camellia-256-ofb & camellia128 & camellia192 & camellia256 \\
cast & cast-cbc & cast5-cbc & cast5-cfb \\
cast5-ecb & cast5-ofb & chacha20 & des \\
des-cbc & des-cfb & des-cfb1 & des-cfb8 \\
des-ecb & des-ede & des-ede-cbc & des-ede-cfb \\
des-ede-ecb & des-ede-ofb & des-ede3 & des-ede3-cbc \\
des-ede3-cfb & des-ede3-cfb1 & des-ede3-cfb8 & des-ede3-ecb \\
des-ede3-ofb & des-ofb & des3 & des3-wrap \\
desx & desx-cbc & id-aes128-wrap & id-aes128-wrap-pad \\
id-aes192-wrap & id-aes192-wrap-pad & id-aes256-wrap & id-aes256-wrap-pad \\
id-smime-alg-CMS3DESwrap & idea & idea-cbc & idea-cfb \\
idea-ecb & idea-ofb & rc2 & rc2-128 \\
rc2-40 & rc2-40-cbc & rc2-64 & rc2-64-cbc \\
rc2-cbc & rc2-cfb & rc2-ecb & rc2-ofb \\
rc4 & rc4-40 & seed & seed-cbc \\
seed-cfb & seed-ecb & seed-ofb & sm4 \\
sm4-cbc & sm4-cfb & sm4-ctr & sm4-ecb \\
sm4-ofb & & & \\
\hline
\end{tabular}
\caption{OpenSSL-supported cipher algorithms}
\end{table}



\chapter{Key Exchange}

Let us do an ephemeral ECDH exchange using OpenSSL.
\begin{enumerate}
\item Generate Alice's EC private key
    \begin{lstlisting}[style=darkbash, language=bash]
    openssl genpkey -algorithm EC -pkeyopt ec_paramgen_curve:P-256 -out alice.pem
    \end{lstlisting} 
\item Generate Bob's EC private key
    \begin{lstlisting}[style=darkbash, language=bash]
    openssl genpkey -algorithm EC -pkeyopt ec_paramgen_curve:P-256 -out bob.pem
    \end{lstlisting} 
\item Extract public keys
    \begin{lstlisting}[style=darkbash, language=bash]
    openssl pkey -in alice.pem -pubout -out alice_pub.pem
    openssl pkey -in bob.pem -pubout -out bob_pub.pem
    \end{lstlisting}
\item Alice derives shared secret using Bob's public key
    \begin{lstlisting}[style=darkbash, language=bash]
    openssl pkeyutl -derive \
  -inkey alice.pem \
  -peerkey bob_pub.pem \
  -out alice_secret.bin
    \end{lstlisting}
\item Bob derives shared secret using Alice's public key
    \begin{lstlisting}[style=darkbash, language=bash]
    openssl pkeyutl -derive \
  -inkey bob.pem \
  -peerkey alice_pub.pem \
  -out bob_secret.bin
    \end{lstlisting}
\item Verify that both sides derived the same secret:
    \begin{lstlisting}[style=darkbash, language=bash]
    sha256sum alice_secret.bin bob_secret.bin
    \end{lstlisting} 
\end{enumerate}

What is happening in principle is
\begin{enumerate}
    \item Alice creates \verb|(alice_private, alice_public)|.
    \item Bob creates (\verb|bob_private, bob_public|).
    \item Alice computes \verb|ECDH(alice_private, bob_public)|.
    \item Bob computes \verb|ECDH(bob_private, alice_public)|.
    \item Both obtain the same shared secret, which can then be fed into a \verb|KDF| to produce encryption keys.
\end{enumerate}

A Key Derivation Function (KDF) is a cryptographic algorithm that turns some input secret material into one or more cryptographic keys.

    










% \chapter{References}

%-- Computation Engineering by Ganesh Gopalakrishnan

\cite{adamson2012set}

% \nocite{*}



\bibliography{../Bibliography/References}
\bibliographystyle{amsalpha}

\printindex




\end{document}